\chapter{Justification}

Although numerous works have attempted to define a function for neurogenesis, the role played by the immature cells it generates
within the circuit they integrate into, as well as what occurs to that circuit during their maturation, remains poorly
understood~\cite{aimoneComputational2016, aimoneModeling2011, berdugo-vegaSharpening2023, faresNeurogenesis2019}.

A number of studies on this topic have focused primarily on the effects of neurogenesis in the DG itself and on the pattern
separation it performs~\cite{berdugo-vegaSharpening2023, kimAdult2024, wangEffect2024}, neglecting its impact on connected
regions such as the CA3 area. This project therefore aims to fill this gap by also examining how iGCs affect the memory
functions of the CA3 area (autoassociation and pattern completion), and how the feedback projections from this area to the DG
influence the system~\cite{myersPattern2011}, with the goal of producing a more holistic and biologically plausible model.

It is also common for a number of works to study the function of a static iGC population~\cite{aimoneComputational2016,
berdugo-vegaSharpening2023}, with models that account for the maturation of these cells over time and the replacement of mGCs
being very rare~\cite{aimoneComputational2009}. Such studies can explain the function (or lack thereof) of these immature cells
only while they remain immature, but are unable to explain how this phase impacts these cells as they mature.

This study is therefore justified by its innovative approach, filling gaps in the existing literature by simulating the
maturation of iGCs and its consequences in CA3.
